\section{Related Work}
\label{sec:barrel-related-work}

Several projects have connected the B family of methods to interactive theorem
provers in order to move proof obligations into a logic with a small,
independently checkable kernel.  An early translation from B to PVS uses type
synthesis to recover the parent-set information needed by the typed target
logic \cite{Bodeveix2002BtoPVS}.  Work in Isabelle/HOL has concentrated on
Event-B: Schmalz formalized its logic, including partial functions and
definitional extensions, in an Isabelle plugin for Rodin
\cite{Schmalz2012RodinIsabelle}, while Ballenghien and Wolff more recently
presented Event-B as a shallow domain-specific language hosted in Isabelle/HOL
\cite{Wolff2024EB_Isabelle}.  These developments either target Event-B and
Rodin or emphasize a formalization of the source logic rather than an
interactive backend for Atelier~B POG files.

Atelier~B obligations have also been exported to automatic provers.  The
multi-prover approach of Mentr\'e et al. and the SMT translation of Deharbe et
al. illustrate this route \cite{mentre2012discharging,deharbe2010automatic}.
The \beer backend developed in this thesis follows the same broad strategy but
adds a verified encoding into SMT-LIB for its supported fragment
\cite{trelat2025beer}.  Such backends are well suited to high-throughput
automatic discharge; their target languages are not intended to become the
interactive proof environment in which a user develops auxiliary mathematics.

\barrel is complementary.  It retains Atelier~B's modelling and
proof-obligation generation stages, consumes the resulting POG, and maps every
supported obligation to an ordinary Lean proposition over Mathlib's typed
sets, relations, orders, and arithmetic.  Its distinctive concern is the
interactive stage: generated statements preserve recognizable B notation,
partial operators carry explicit well-definedness evidence, and successful
proofs become named, kernel-checked Lean theorems.  The goal is therefore not
to formalize B anew, nor to replace automatic backends, but to make Lean and
its library available as a modern second prover for an existing Atelier~B
development.
